%==============================================================================
% Assignment 1 - Generative AI - technical report
%
% This is a SKELETON. Every \PLACEHOLDER command below marks something only you
% can fill in: a measured number, an observation from your own runs, a decision
% you made and can defend in the viva. Compile it once before you start writing
% so you know the toolchain works:
%
%   pdflatex main && bibtex main && pdflatex main && pdflatex main
%
% Search for "PLACEHOLDER" before submitting - none should remain.
%==============================================================================
\documentclass[conference]{IEEEtran}
\IEEEoverridecommandlockouts

\usepackage{cite}
\usepackage{amsmath,amssymb,amsfonts}
\usepackage{algorithmic}
\usepackage{graphicx}
\usepackage{textcomp}
\usepackage{xcolor}
\usepackage{booktabs}
\usepackage{multirow}
\usepackage{url}
\usepackage[hidelinks]{hyperref}

\def\BibTeX{{\rm B\kern-.05em{\sc i\kern-.025em b}\kern-.08em T\kern-.1667em\lower.7ex\hbox{E}\kern-.125emX}}

% Makes unfinished text impossible to miss in the compiled PDF.
\newcommand{\PLACEHOLDER}[1]{\textcolor{red}{\textbf{[FILL IN: #1]}}}

\begin{document}

\title{Universal, Hard-Routed and Soft Mixture-of-Experts Image Restoration\\
with Style-Conditioned Face-to-Sketch Generation}

\author{\IEEEauthorblockN{\PLACEHOLDER{your name}}
\IEEEauthorblockA{\textit{\PLACEHOLDER{department}} \\
\textit{\PLACEHOLDER{university}}\\
\PLACEHOLDER{city, country} \\
\PLACEHOLDER{email}}
}

\maketitle

\begin{abstract}
\PLACEHOLDER{150--200 words. State the problem, name the four systems, give
your two or three headline numbers (for example the universal model's mean test
PSNR, the classifier's macro F1, the SSIM gap between oracle and predicted
routing, and the cGAN's test L1), and close with one sentence on the deployed
application. Write this last, once the numbers exist.}
\end{abstract}

\begin{IEEEkeywords}
denoising autoencoder, image restoration, mixture of experts, conditional GAN,
hyperparameter optimisation, ONNX
\end{IEEEkeywords}

%==============================================================================
\section{Introduction}
%==============================================================================
Image restoration systems are usually built for one degradation at a time. A
denoiser trained on Gaussian noise does little for a blurred photograph, and an
inpainting model has no reason to remove impulse noise. This report studies
three ways of handling several degradations with one system, and a fourth,
separate problem in paired image-to-image translation.

\PLACEHOLDER{Two or three paragraphs: why multi-corruption restoration is hard,
what a single shared representation has to trade off, and why routing (hard or
soft) is an alternative worth measuring. Close with an explicit contribution
list.}

The contributions of this work are:
\begin{itemize}
  \item a single universal denoising autoencoder covering clean input,
        impulse noise, Gaussian blur and rectangular occlusion;
  \item a hard-routed system pairing a four-class corruption classifier with
        three specialist autoencoders, measured under both oracle and predicted
        routing;
  \item a differentiable soft mixture-of-experts restorer, initialised from the
        hard-routed components and fine-tuned jointly;
  \item a style-conditioned conditional GAN for face-to-sketch generation; and
  \item a containerised application exposing all four systems through one
        browser interface backed by ONNX Runtime.
\end{itemize}

%==============================================================================
\section{Related Work}
%==============================================================================
\textbf{Denoising autoencoders.} Vincent et al.~\cite{vincent2008} introduced
corruption as a training signal rather than a nuisance, showing that
reconstructing a clean target from a damaged input yields robust
representations. \PLACEHOLDER{Add the two or three restoration papers you
actually read and say what you took from each.}

\textbf{Mixtures of experts.} Adaptive mixtures of local
experts~\cite{jacobs1991} learn a gate that assigns inputs to specialised
sub-networks. Sparsely-gated variants~\cite{shazeer2017} and Switch
routing~\cite{fedus2022} scale the idea but reintroduce a discrete decision;
soft mixtures~\cite{puigcerver2023} keep the combination differentiable.
\PLACEHOLDER{Explain which of these informed your Task 3 design and how the
load-balancing term in your loss relates to theirs.}

\textbf{Paired image-to-image translation.} pix2pix~\cite{isola2017} combines a
U-Net generator~\cite{ronneberger2015} with a PatchGAN discriminator and an
$L_1$ reconstruction term; the conditional GAN formulation itself is due to
Mirza and Osindero~\cite{mirza2014}. \PLACEHOLDER{Add the facial-sketch
synthesis work you consulted and say how your style conditioning differs.}

\textbf{Quality metrics.} SSIM~\cite{wang2004} is used throughout as a
structural complement to pixel error. \PLACEHOLDER{If you considered LPIPS or
FID and rejected them, say why here.}

%==============================================================================
\section{Dataset Preparation}
%==============================================================================
\subsection{Oxford-IIIT Pet (Tasks 1--3)}
The official \texttt{trainval} collection~\cite{parkhi2012} is split 80/20 into
training and validation partitions with \texttt{numpy.random.RandomState(42)}.
The official test collection is untouched until final evaluation. All images
are converted to RGB and resized to $128\times128$. The identical split is
reused by Tasks 1, 2 and 3.

Sizes: \PLACEHOLDER{train / validation / test image counts from your run}.

\subsection{Corruption Model}
Corruptions are generated inside the data loader, so an image is damaged
differently every time it is read and no corrupted copies are stored. Each
training sample draws one of the four conditions in
Table~\ref{tab:corruptions} with equal probability, and the drawn class is the
label used later to train the classifier and the gate.

\begin{table}[t]
\caption{Corruption configuration}
\label{tab:corruptions}
\centering
\begin{tabular}{@{}lp{5.4cm}@{}}
\toprule
\textbf{Condition} & \textbf{Training configuration} \\
\midrule
Clean & Image passed through unchanged. \\
Salt \& pepper & $p \sim \mathcal{U}(0.02, 0.15)$; selected pixels set to
  black or white with equal probability. \\
Gaussian blur & kernel $\in \{3,5,7\}$, $\sigma \sim \mathcal{U}(0.5, 2.5)$. \\
Occlusion & 1--3 black rectangles jointly covering 10--35\% of the area, placed
  at random. \\
\bottomrule
\end{tabular}
\end{table}

Validation and test corruptions are deterministic. A manifest generated once
stores the corruption type, severity, blur settings, rectangle coordinates and
seed for every entry, so every evaluation in this report runs on byte-identical
inputs. Test severities are fixed at $p \in \{0.03, 0.08, 0.15\}$ for impulse
noise, $(k,\sigma) \in \{(3,0.7),(5,1.5),(7,2.5)\}$ for blur, and nominal
coverages of 10\%, 20\% and 35\% using one, two and three rectangles.

Because rectangles may overlap, realised occlusion coverage is slightly below
the nominal target; the manifest records the measured value. Mean realised
coverage: \PLACEHOLDER{report the three measured values from
\texttt{test\_manifest.json}}.

\subsection{FS2K (Task 4)}
FS2K~\cite{fan2022} provides 2{,}104 photograph--sketch pairs in three sketch
styles. The official training portion is split 85/15 with seed 42, stratified by
style; the official test portion is reserved. Both images of a pair receive
\emph{identical} spatial augmentation, since independent transforms would
destroy the pixel correspondence the $L_1$ term depends on.

Per-style counts: \PLACEHOLDER{train / validation / test counts per style}.

%==============================================================================
\section{Task 1: Universal Denoising Autoencoder}
%==============================================================================
\subsection{Architecture}
The encoder halves the spatial resolution four times while doubling channel
width, producing a latent tensor of $d \times 8 \times 8$; the decoder mirrors
it with nearest-neighbour upsampling followed by convolution, avoiding the
checkerboard artefacts of transposed convolutions~\cite{odena2016}. With the
selected $d$ the compression ratio is \PLACEHOLDER{ratio reported by
\texttt{compression\_ratio()}}, so the bottleneck is genuine rather than a wide
identity path.

\PLACEHOLDER{If you enabled the coarse skip connection, justify it here with a
measured comparison: PSNR/SSIM with and without, and what the error maps show.
If you left it off, say that and give the reason.}

\begin{figure}[t]
\centering
% \includegraphics[width=\columnwidth]{figures/task1_architecture.pdf}
\caption{Universal autoencoder architecture. \PLACEHOLDER{draw this}}
\label{fig:task1arch}
\end{figure}

\subsection{Loss}
Training minimises
\begin{equation}
\mathcal{L}_{\text{rec}} = \alpha \mathcal{L}_{1}(x, \hat{x})
  + (1-\alpha)\bigl(1 - \mathrm{SSIM}(x, \hat{x})\bigr),
\end{equation}
where $x$ is the clean target and $\hat{x}=D(E(\tilde{x}))$ the reconstruction
of the corrupted input $\tilde{x}$. SSIM uses an $11\times11$ Gaussian window
with $\sigma=1.5$ and the standard stabilising
constants~\cite{wang2004}. Rather than accepting $\alpha=0.8$, the value is
selected by the Optuna study; the selected value is \PLACEHOLDER{$\alpha$}.

\subsection{Hyperparameter Optimisation}
\begin{table}[t]
\caption{Task 1 Optuna search space and selected values}
\label{tab:task1optuna}
\centering
\begin{tabular}{@{}lll@{}}
\toprule
\textbf{Parameter} & \textbf{Range} & \textbf{Selected} \\
\midrule
Learning rate & $[10^{-4}, 5\times10^{-3}]$ log & \PLACEHOLDER{} \\
Batch size & $\{16,32,64,128\}$ & \PLACEHOLDER{} \\
Bottleneck dimension & $\{32,64,128,256\}$ & \PLACEHOLDER{} \\
Base channels & $\{16,32,48,64\}$ & \PLACEHOLDER{} \\
Dropout & $[0.0, 0.4]$ & \PLACEHOLDER{} \\
Loss weight $\alpha$ & $[0.3, 0.95]$ & \PLACEHOLDER{} \\
Weight decay & $[10^{-6}, 10^{-3}]$ log & \PLACEHOLDER{} \\
\bottomrule
\end{tabular}
\end{table}

The objective maximises $\mathrm{SSIM} + \mathrm{PSNR}/50$ on the validation
manifest, a single scalar that refuses configurations which win on pixel error
by over-smoothing. \PLACEHOLDER{Report trials completed, trials pruned, best
trial number and value, and which parameters mattered according to the
importance plot.}

\subsection{Results}
\begin{table}[t]
\caption{Task 1 test performance by corruption and severity}
\label{tab:task1results}
\centering
\begin{tabular}{@{}llcc@{}}
\toprule
\textbf{Condition} & \textbf{Severity} & \textbf{PSNR (dB)} & \textbf{SSIM} \\
\midrule
Clean & --- & \PLACEHOLDER{} & \PLACEHOLDER{} \\
\midrule
\multirow{3}{*}{Salt \& pepper} & low & \PLACEHOLDER{} & \PLACEHOLDER{} \\
 & medium & \PLACEHOLDER{} & \PLACEHOLDER{} \\
 & high & \PLACEHOLDER{} & \PLACEHOLDER{} \\
\midrule
\multirow{3}{*}{Blur} & low & \PLACEHOLDER{} & \PLACEHOLDER{} \\
 & medium & \PLACEHOLDER{} & \PLACEHOLDER{} \\
 & high & \PLACEHOLDER{} & \PLACEHOLDER{} \\
\midrule
\multirow{3}{*}{Occlusion} & low & \PLACEHOLDER{} & \PLACEHOLDER{} \\
 & medium & \PLACEHOLDER{} & \PLACEHOLDER{} \\
 & high & \PLACEHOLDER{} & \PLACEHOLDER{} \\
\bottomrule
\end{tabular}
\end{table}

\PLACEHOLDER{Interpret the table. Which corruption is hardest and why? A useful
observation is that impulse noise is locally recoverable while occlusion
requires the decoder to hallucinate content it never saw, so SSIM degrades much
faster with occlusion severity than with noise severity --- but only claim this
if your numbers show it.}

Figure~\ref{fig:task1samples} shows twelve representative reconstructions with
their absolute error maps.

\begin{figure}[t]
\centering
% \includegraphics[width=\columnwidth]{figures/task1_representative.png}
\caption{Clean target, corrupted input, reconstruction and absolute error for
twelve representative test images.}
\label{fig:task1samples}
\end{figure}

\subsubsection{Failure cases}
\PLACEHOLDER{Discuss four specific failures from
\texttt{task1\_failure\_cases.json}. For each: what the input was, what the
model produced, and the mechanism you think is responsible. Generic statements
such as ``the model struggled'' are not sufficient --- point at the error map
and say what it shows.}

%==============================================================================
\section{Task 2: Classification and Hard-Routed Specialists}
%==============================================================================
\subsection{Corruption Classifier}
The classifier predicts $p = f(\tilde{x}) = [p_{\text{clean}}, p_{\text{salt}},
p_{\text{blur}}, p_{\text{occ}}]$ and is trained with multiclass cross-entropy
on class-balanced batches, obtained by indexing the dataset as (image,
corruption) pairs so that every image appears once per class in an epoch.

\begin{table}[t]
\caption{Classifier performance on the test manifest}
\label{tab:task2cls}
\centering
\begin{tabular}{@{}lccc@{}}
\toprule
\textbf{Class} & \textbf{Precision} & \textbf{Recall} & \textbf{F1} \\
\midrule
Clean & \PLACEHOLDER{} & \PLACEHOLDER{} & \PLACEHOLDER{} \\
Salt \& pepper & \PLACEHOLDER{} & \PLACEHOLDER{} & \PLACEHOLDER{} \\
Blur & \PLACEHOLDER{} & \PLACEHOLDER{} & \PLACEHOLDER{} \\
Occlusion & \PLACEHOLDER{} & \PLACEHOLDER{} & \PLACEHOLDER{} \\
\midrule
Macro average & \PLACEHOLDER{} & \PLACEHOLDER{} & \PLACEHOLDER{} \\
\bottomrule
\end{tabular}
\end{table}

\begin{figure}[t]
\centering
% \includegraphics[width=0.85\columnwidth]{figures/task2_confusion_matrix.png}
\caption{Normalised four-class confusion matrix.}
\label{fig:task2cm}
\end{figure}

\PLACEHOLDER{Read the confusion matrix out loud in prose. Which pair is
confused, at which severity, and why? Mild blur against clean, and very light
impulse noise against clean, are the plausible candidates --- check whether your
matrix agrees before asserting it.}

\subsection{Specialists and Routing}
Three specialists are trained, each on one corruption, sharing an architecture
found by a single Optuna study but with independently trained weights. At
inference the predicted class selects the branch, and a clean prediction
bypasses restoration entirely:
\begin{equation}
\hat{x} = \begin{cases}
\tilde{x} & c = \text{clean}\\
E_{\text{salt}}(\tilde{x}) & c = \text{salt}\\
E_{\text{blur}}(\tilde{x}) & c = \text{blur}\\
E_{\text{occ}}(\tilde{x}) & c = \text{occlusion}
\end{cases}
\end{equation}

\begin{table}[t]
\caption{Oracle versus predicted routing}
\label{tab:task2routing}
\centering
\begin{tabular}{@{}lcccc@{}}
\toprule
\multirow{2}{*}{\textbf{Condition}} & \multicolumn{2}{c}{\textbf{Oracle}} &
\multicolumn{2}{c}{\textbf{Predicted}} \\
\cmidrule(lr){2-3}\cmidrule(lr){4-5}
 & PSNR & SSIM & PSNR & SSIM \\
\midrule
Clean & \PLACEHOLDER{} & \PLACEHOLDER{} & \PLACEHOLDER{} & \PLACEHOLDER{} \\
Salt \& pepper & \PLACEHOLDER{} & \PLACEHOLDER{} & \PLACEHOLDER{} & \PLACEHOLDER{} \\
Blur & \PLACEHOLDER{} & \PLACEHOLDER{} & \PLACEHOLDER{} & \PLACEHOLDER{} \\
Occlusion & \PLACEHOLDER{} & \PLACEHOLDER{} & \PLACEHOLDER{} & \PLACEHOLDER{} \\
\midrule
Overall & \PLACEHOLDER{} & \PLACEHOLDER{} & \PLACEHOLDER{} & \PLACEHOLDER{} \\
\bottomrule
\end{tabular}
\end{table}

The oracle column isolates specialist quality; the predicted column is the
operational system. The gap between them is the price of classifier error.
\PLACEHOLDER{Quantify that gap and then discuss the specific misrouting cases
in \texttt{task2\_test\_metrics.json}: which misclassifications were harmless
(two experts that do similar things) and which were destructive (sending an
occluded image to the deblurring expert)?}

\PLACEHOLDER{Compare Task 2 against Task 1 directly. Did specialisation beat
the universal model, and on which corruptions? If the universal model won
anywhere, say so and explain it rather than hiding it.}

%==============================================================================
\section{Task 3: Soft Mixture of Experts}
%==============================================================================
\subsection{Formulation}
The gate produces continuous weights over four branches, the first being an
identity path for clean input:
\begin{equation}
w = \mathrm{softmax}\!\left(\frac{g(\tilde{x})}{T}\right),\quad
\hat{x} = w_0\tilde{x} + w_1 E_{\text{salt}}(\tilde{x})
  + w_2 E_{\text{blur}}(\tilde{x}) + w_3 E_{\text{occ}}(\tilde{x}).
\end{equation}
Because the combination is differentiable, gate and experts train through the
reconstruction error. The gate is initialised from the Task 2 classifier and
the experts from the Task 2 specialists; a warm-up stage trains the gate alone
with the experts frozen, after which everything is fine-tuned at a smaller
learning rate.

The joint objective is
\begin{equation}
\mathcal{L} = \lambda_1 \mathcal{L}_1 + \lambda_2 (1-\mathrm{SSIM})
  + \lambda_3 \mathcal{L}_{\text{CE}} + \lambda_4 \mathcal{L}_{\text{balance}},
\end{equation}
with
$\mathcal{L}_{\text{balance}} = \sum_{i=1}^{4} (\bar{w}_i - \tfrac{1}{4})^2$
computed over class-balanced batches, where $\bar{w}_i$ is the mean weight of
branch $i$. \PLACEHOLDER{Justify this regulariser against the entropy
alternative. The uniform target is only correct because the batches are class
balanced --- say so, and say what you would use instead if they were not.}

\subsection{Routing Behaviour}
\begin{table}[t]
\caption{Mean routing weight by true corruption}
\label{tab:task3routing}
\centering
\begin{tabular}{@{}lcccc@{}}
\toprule
\textbf{True input} & \textbf{Identity} & \textbf{Salt} & \textbf{Blur} &
\textbf{Occlusion} \\
\midrule
Clean & \PLACEHOLDER{} & \PLACEHOLDER{} & \PLACEHOLDER{} & \PLACEHOLDER{} \\
Salt \& pepper & \PLACEHOLDER{} & \PLACEHOLDER{} & \PLACEHOLDER{} & \PLACEHOLDER{} \\
Blur & \PLACEHOLDER{} & \PLACEHOLDER{} & \PLACEHOLDER{} & \PLACEHOLDER{} \\
Occlusion & \PLACEHOLDER{} & \PLACEHOLDER{} & \PLACEHOLDER{} & \PLACEHOLDER{} \\
\bottomrule
\end{tabular}
\end{table}

\begin{figure}[t]
\centering
% \includegraphics[width=0.9\columnwidth]{figures/task3_routing_heatmap.png}
\caption{Routing weights by true corruption and severity.}
\label{fig:task3heat}
\end{figure}

\PLACEHOLDER{Answer the three questions explicitly: (i) show an example where a
single expert dominates; (ii) show an example where the weight is genuinely
split, and say whether the blended output is better than either branch alone;
(iii) state whether any expert became inactive or captured inputs unrelated to
its specialisation. \texttt{task3\_test\_metrics.json} lists blended cases and
flags inactive branches.}

\PLACEHOLDER{Report the effect of temperature $T$. Small $T$ approaches hard
routing, large $T$ spreads weight. Give the selected value and what the search
showed about it.}

\subsection{Cost}
Soft routing evaluates every expert on every input, so inference cost is
roughly three times the hard-routed system.
\PLACEHOLDER{Give measured per-image inference times for Tasks 1, 2 and 3 from
the application, and state whether the accuracy gain justifies the cost.}

%==============================================================================
\section{Task 4: Style-Conditioned Face-to-Sketch cGAN}
%==============================================================================
\subsection{Architecture}
The generator is a U-Net with six downsampling and six upsampling stages; the
discriminator is a PatchGAN classifying overlapping local
regions~\cite{isola2017}. The style condition is an \texttt{nn.Embedding} over
the three FS2K categories, broadcast across the spatial grid and concatenated
to the generator input, and additionally projected into the bottleneck. The
discriminator receives the same embedding, so a sketch paired with the wrong
style condition is detectable as fake --- without this the generator is free to
ignore the condition entirely.

\subsection{Objective}
\begin{equation}
\mathcal{L}_G = \mathcal{L}_{\text{adv}}
  + \lambda_{L1}\,\mathcal{L}_{1}\bigl(y, G(x, s)\bigr)
\end{equation}
Adversarial terms use binary cross-entropy with logits and one-sided label
smoothing on the real target. The conventional $\lambda_{L1}=100$ is used as a
starting point only; the selected value is \PLACEHOLDER{$\lambda_{L1}$}.

\subsection{Results}
\begin{table}[t]
\caption{Task 4 test performance by style}
\label{tab:task4}
\centering
\begin{tabular}{@{}lcccc@{}}
\toprule
\textbf{Style} & \textbf{$L_1$} & \textbf{SSIM} & \textbf{PSNR} & \textbf{n} \\
\midrule
Style 1 & \PLACEHOLDER{} & \PLACEHOLDER{} & \PLACEHOLDER{} & \PLACEHOLDER{} \\
Style 2 & \PLACEHOLDER{} & \PLACEHOLDER{} & \PLACEHOLDER{} & \PLACEHOLDER{} \\
Style 3 & \PLACEHOLDER{} & \PLACEHOLDER{} & \PLACEHOLDER{} & \PLACEHOLDER{} \\
Overall & \PLACEHOLDER{} & \PLACEHOLDER{} & \PLACEHOLDER{} & \PLACEHOLDER{} \\
\bottomrule
\end{tabular}
\end{table}

\begin{figure}[t]
\centering
% \includegraphics[width=\columnwidth]{figures/task4_per_style.png}
\caption{Photograph, ground-truth sketch and generated sketch for each style.}
\label{fig:task4grid}
\end{figure}

\subsubsection{Does the condition do anything?}
\texttt{evaluate\_cgan.py} generates the same photograph under all three
conditions and reports the mean absolute difference between outputs.
Measured differences: \PLACEHOLDER{values}.
\PLACEHOLDER{If the differences are near zero the generator is ignoring the
style embedding, which is the most common failure of this architecture. Report
it honestly and explain what you tried.}

\subsubsection{Training stability}
\PLACEHOLDER{Discuss the four logged curves. A discriminator loss collapsing
toward zero while the generator adversarial loss climbs indicates the
discriminator has won; describe what you observed and any intervention
(learning-rate ratio, label smoothing, fewer discriminator steps) you made.}

\subsubsection{Failure cases}
\PLACEHOLDER{Four specific failures: unusual poses, glasses, hair boundaries,
lighting. Show them.}

%==============================================================================
\section{Application and Deployment}
%==============================================================================
\subsection{Interface Design}
The layout was designed in Google Stitch before implementation.
\PLACEHOLDER{Insert the Stitch screenshots and say what changed between the
design and the built interface, and why.}

\begin{figure}[t]
\centering
% \includegraphics[width=\columnwidth]{figures/stitch_design.png}
\caption{Original Google Stitch design.}
\label{fig:stitch}
\end{figure}

\subsection{Architecture}
A React and Tailwind CSS frontend communicates with a FastAPI backend over five
endpoints: health, universal restoration, hard routing, soft mixture and
face-to-sketch. The backend validates the upload, resizes to $128\times128$,
runs ONNX Runtime and returns the result with routing information and timing.
Both services run as Docker containers started by one Compose command; nginx
serves the built frontend and proxies \texttt{/api} to the backend, so the
browser bundle contains no hard-coded backend address.

\subsection{ONNX Export and Verification}
All seven inference models are exported at opset 17. The Task 3 export contains
the entire pipeline --- gate, temperature-scaled softmax, three experts and the
weighted sum --- as a single graph, so the served model cannot drift from the
trained one. Agreement with PyTorch was verified on random inputs; the largest
observed deviation was \PLACEHOLDER{value} against a tolerance of $10^{-4}$.

\begin{figure}[t]
\centering
% \includegraphics[width=\columnwidth]{figures/app_screenshots.png}
\caption{The four workspaces.}
\label{fig:app}
\end{figure}

%==============================================================================
\section{Limitations}
%==============================================================================
\PLACEHOLDER{Be specific and honest. Candidates: single corruption per image at
training time, so mixed degradations are out of distribution; fixed
$128\times128$ resolution; the occlusion expert can only hallucinate plausible
texture, not recover content; soft routing costs three expert evaluations per
image; FS2K is small and demographically narrow, so the generator's behaviour
on faces unlike the training distribution is unmeasured; CPU-only inference in
the container.}

%==============================================================================
\section{Conclusion}
%==============================================================================
\PLACEHOLDER{Three or four sentences tying the measured results back to the
question the assignment poses: when is one shared model enough, when does
routing pay for itself, and what did the soft mixture buy over the hard one?}

%==============================================================================
\appendices
\section{Repository and Demonstration}
Source code: \url{PLACEHOLDER-github-url}\\
Trained models: \url{PLACEHOLDER-release-or-lfs-url}\\
Demonstration video: \url{PLACEHOLDER-youtube-url}

\section{Use of AI Assistance}
\PLACEHOLDER{Required appendix. For each tool, state what it was used for and
how its output was tested or corrected. Be concrete --- name the files, name the
bugs you found. For example: a coding assistant produced the initial data
pipeline and training loops; the corruption statistics were verified against the
specification with \texttt{scripts/smoke\_test.py}, which checks that the
realised salt-and-pepper fraction matches the sampled probability and that
occlusion coverage matches the nominal target; ONNX agreement with PyTorch was
verified with \texttt{src/export/verify\_onnx.py}. List anything the assistant
got wrong and how you fixed it --- that is the part the evaluator is looking
for.}

\begin{thebibliography}{00}
\bibitem{vincent2008} P. Vincent, H. Larochelle, Y. Bengio, and P.-A. Manzagol,
``Extracting and composing robust features with denoising autoencoders,'' in
\emph{Proc. ICML}, 2008, pp. 1096--1103.

\bibitem{wang2004} Z. Wang, A. C. Bovik, H. R. Sheikh, and E. P. Simoncelli,
``Image quality assessment: From error visibility to structural similarity,''
\emph{IEEE Trans. Image Process.}, vol. 13, no. 4, pp. 600--612, 2004.

\bibitem{jacobs1991} R. A. Jacobs, M. I. Jordan, S. J. Nowlan, and G. E. Hinton,
``Adaptive mixtures of local experts,'' \emph{Neural Computation}, vol. 3,
no. 1, pp. 79--87, 1991.

\bibitem{shazeer2017} N. Shazeer et al., ``Outrageously large neural networks:
The sparsely-gated mixture-of-experts layer,'' in \emph{Proc. ICLR}, 2017.

\bibitem{fedus2022} W. Fedus, B. Zoph, and N. Shazeer, ``Switch transformers:
Scaling to trillion parameter models with simple and efficient sparsity,''
\emph{J. Mach. Learn. Res.}, vol. 23, no. 120, pp. 1--39, 2022.

\bibitem{puigcerver2023} J. Puigcerver, C. Riquelme, B. Mustafa, and N. Houlsby,
``From sparse to soft mixtures of experts,'' \emph{arXiv:2308.00951}, 2023.

\bibitem{isola2017} P. Isola, J.-Y. Zhu, T. Zhou, and A. A. Efros,
``Image-to-image translation with conditional adversarial networks,'' in
\emph{Proc. CVPR}, 2017, pp. 1125--1134.

\bibitem{ronneberger2015} O. Ronneberger, P. Fischer, and T. Brox, ``U-Net:
Convolutional networks for biomedical image segmentation,'' in \emph{Proc.
MICCAI}, 2015, pp. 234--241.

\bibitem{mirza2014} M. Mirza and S. Osindero, ``Conditional generative
adversarial nets,'' \emph{arXiv:1411.1784}, 2014.

\bibitem{goodfellow2014} I. Goodfellow et al., ``Generative adversarial nets,''
in \emph{Proc. NeurIPS}, 2014, pp. 2672--2680.

\bibitem{parkhi2012} O. M. Parkhi, A. Vedaldi, A. Zisserman, and C. V. Jawahar,
``Cats and dogs,'' in \emph{Proc. CVPR}, 2012, pp. 3498--3505.

\bibitem{fan2022} D.-P. Fan et al., ``Facial-sketch synthesis: A new
challenge,'' \emph{Machine Intelligence Research}, vol. 19, no. 4,
pp. 257--287, 2022.

\bibitem{akiba2019} T. Akiba, S. Sano, T. Yanase, T. Ohta, and M. Koyama,
``Optuna: A next-generation hyperparameter optimization framework,'' in
\emph{Proc. KDD}, 2019, pp. 2623--2631.

\bibitem{odena2016} A. Odena, V. Dumoulin, and C. Olah, ``Deconvolution and
checkerboard artifacts,'' \emph{Distill}, 2016.

\bibitem{kingma2015} D. P. Kingma and J. Ba, ``Adam: A method for stochastic
optimization,'' in \emph{Proc. ICLR}, 2015.

\bibitem{loshchilov2019} I. Loshchilov and F. Hutter, ``Decoupled weight decay
regularization,'' in \emph{Proc. ICLR}, 2019.
\end{thebibliography}

\end{document}
